\chapter{The 2$\times$2 frame}
\label{ch:quadrant}

\newthought{If you have read this far} you have a list of architectures
that all claim to do something interesting with long context: dense
transformers, FlashAttention, MQA/GQA, sliding windows, Longformer,
BigBird, S4, Mamba, Hyena, Performer, RetNet, Jamba, SAMBA, Griffin,
Subquadratic. The list keeps growing. Without a frame to compare them
they look like a zoo.

This chapter is the frame. Two axes, four quadrants, every architecture
gets a pin.

\section{The two axes}

\paragraph{Axis 1 — Routing.} How does the model choose which past
information to use at each step?

\begin{description}
\item[Position-fixed routing.] The set of relevant past tokens is
determined entirely by token \emph{positions}, not by their content.
``Attend to the last \(w\) tokens.'' ``Attend to every \(k\)-th
position.'' ``Convolve with a fixed kernel.'' ``Maintain a fixed-size
running state.'' Cheap and parallelisable. Bad at queries that need a
specific fact at an unpredictable position.
\item[Content-dependent routing.] The model uses the current query to
\emph{select} which past tokens matter, based on a similarity score
between the query and the past tokens' representations. Soft selection
(\emph{softmax over all positions}) is what dense attention does. Hard
selection (\emph{pick top-\(k\)}) is what content-dependent sparse
attention does.
\end{description}

\paragraph{Axis 2 — Scaling.} What is the asymptotic compute cost as a
function of context length \(n\)?

\begin{description}
\item[Quadratic.] \(\bigO(n^2)\). Every pair of tokens scored.
\item[Linear (or near-linear).] \(\bigO(n)\) or \(\bigO(n \log n)\).
The total compute grows in step with the input.
\end{description}

\section{The four quadrants}

\begin{figure}
\centering
\begin{tikzpicture}[
  every node/.style={font=\small},
  qbox/.style={draw=gray!50, text width=3.6cm, align=center,
               minimum height=3.0cm, rounded corners, inner sep=4pt},
  axislabel/.style={font=\footnotesize}
]
  \node[qbox] (tl) at (0,0)   {\textbf{Position-fixed, quadratic}\\[2pt]\textit{(empty for good reason)}};
  \node[qbox] (tr) at (4.6,0) {\textbf{Content-dependent, quadratic}\\[2pt]Transformer\\FlashAttention\\DeepSeek SA*};
  \node[qbox] (bl) at (0,-3.6){\textbf{Position-fixed, linear}\\[2pt]Sliding window\\Longformer / BigBird\\S4 / Mamba (lossy)\\Hyena (long conv)};
  \node[qbox] (br) at (4.6,-3.6){\textbf{Content-dependent, linear}\\[2pt]\textbf{SSA / SubQ}\\(content-routed\\sparse mechanisms)};

  \draw[-{Latex[length=2mm]}, thick] (-3.2, -1.8) -- (-2.0, -1.8);
  \node[axislabel] at (-2.6, -1.5) {routing};
  \node[axislabel] at (-3.2, -2.2) {pos.\ fixed};
  \node[axislabel] at (-2.0, -2.2) {content};
  \draw[-{Latex[length=2mm]}, thick] (-3.2, -2.6) -- (-3.2, -1.2);
  \node[axislabel, rotate=90] at (-3.6, -2.6) {linear};
  \node[axislabel, rotate=90] at (-3.6, -1.2) {quadratic};
\end{tikzpicture}
\caption{The 2$\times$2 frame. Top-right is dense attention and its
implementation-only optimisations. Bottom-left is the position-fixed
linear-cost zoo. Bottom-right — content-dependent and linear — is the
quadrant SSA aims at. Top-left is empty because spending \(\bigO(n^2)\)
on a content-blind mechanism is strictly dominated.}
\label{fig:quadrant}
\end{figure}

\paragraph{Top-left (position-fixed, quadratic).} Empty. Spending \(n^2\)
compute without using content to choose which pairs to score is strictly
worse than dense attention; nobody builds these.

\paragraph{Top-right (content-dependent, quadratic).} The classic
transformer. FlashAttention, MQA, GQA, paged caches, RoPE/YaRN — every
mitigation in Part~IV — lives here. The asymptote is \(n^2\); the
constants are spectacular by 2026 standards. DeepSeek's recently published
DSA \citep{deepseek2024dsa} is also content-dependent but applies sparsity
at training-time-decided sparsity patterns, blurring the line.

\paragraph{Bottom-left (position-fixed, linear).} The two big sub-families
of Part~V live here: fixed-pattern sparse (Chapter~\ref{ch:sparse}) and
state-space / convolutional models (Chapters~\ref{ch:ssm}--\ref{ch:linattn}).
They are linear in \(n\) but content-blind: routing is decided by
position or by a fixed kernel, not by the data. This is the dominant
flavour of ``efficient'' attention and it has known weak spots on
associative-recall tasks.

\paragraph{Bottom-right (content-dependent, linear).} The hard quadrant.
You need to keep the dot-product-routing virtue of dense attention, while
limiting the number of (query, key) pairs you actually score to be
\emph{linear} in \(n\). This requires a non-trivial selection mechanism:
given a query, find the small set of keys that matter, in $\bigO(\log n)$
or $\bigO(n)$ work, without doing the full $n \times n$ dot product first.

\keyidea{The interesting frontier in long-context architectures is the
bottom-right quadrant: content-dependent routing at linear cost. Every
major architecture in Part~V occupies a different quadrant; SSA / SubQ
explicitly aims for the bottom-right.}

\fillme{FILL-16-01}{Sidebar: where does DeepSeek Sparse Attention sit?}{%
  ${\sim}300$-word sidebar placing DeepSeek Sparse Attention (DSA) on the
  quadrant. Argue that DSA is content-dependent (selection is data-driven
  via a learned router) but its computational cost remains roughly
  $\bigO(n^2)$ in the worst case because the routing structure is denser
  than e.g.\ SSA's. Cite DeepSeek's 2024 technical report.}{%
  DeepSeek 2024 technical report; the company's blog post.}{%
  ${\sim}300$ words; no equations.}

\fillme{FILL-16-02}{Worked example: cost curves overlaid for \(n \in \{10^3, 10^4, 10^5, 10^6\}\)}{%
  Produce a single \texttt{pgfplots} figure with four curves on the same
  log-log axes: dense \(\bigO(n^2)\), sliding-window \(\bigO(n w)\) at
  \(w = 4096\), Mamba \(\bigO(n)\), and SSA \(\bigO(n \log n)\).
  Normalise so all four curves coincide at \(n = 10^3\). Caption: this is
  the picture that motivates the entire bottom row of the quadrant.}{%
  Asymptotics from Chapters~\ref{ch:why-quadratic}, \ref{ch:sparse},
  \ref{ch:ssm}, \ref{ch:ssa}.}{%
  One small \texttt{pgfplots} figure + 3 sentences caption.}
